\section{Исходный код}
Программа \texttt{main.py} написана на Python 3 и использует библиотеки
\texttt{requests} и \texttt{BeautifulSoup}. Основные функции: \texttt{api} (запрос
к MediaWiki API с повторами), \texttt{random\_pages} (список случайных статей),
\texttt{extract\_text} (очистка HTML), \texttt{download} (загрузка и сохранение
одной статьи) и \texttt{crawl} (основной цикл).

Этапы работы:
\begin{enumerate}
    \item получение пачки из 500 случайных статей (без перенаправлений);
    \item отбрасывание уже скачанных идентификаторов;
    \item параллельная (в нескольких потоках) загрузка HTML каждой статьи;
    \item выделение текста, пропуск коротких статей;
    \item сохранение HTML-документа и JSON-метаданных.
\end{enumerate}

Учтены особенности реального API: временные сбои (повторные попытки), ограничение
частоты запросов (код 429 и заголовок \texttt{Retry-After}), пропуск удалённых и
недоступных статей, возможность остановить и продолжить загрузку.

Запрос к API с повторами и обработкой ограничения частоты:

\begin{lstlisting}[language=Python]
def api(**params):
    for attempt in range(5):
        try:
            r = session.get(API, params={"format": "json",
                            "formatversion": 2, **params}, timeout=30)
            if r.status_code == 200:
                return r.json()
            if r.status_code == 429:   # rate limited
                time.sleep(int(r.headers.get("Retry-After", 10)))
                continue
        except (requests.RequestException, ValueError):
            pass
        time.sleep(5 * (attempt + 1))
    raise RuntimeError(f"API failed: {params}")
\end{lstlisting}

Выделение текста из HTML статьи:

\begin{lstlisting}[language=Python]
def extract_text(html):
    soup = BeautifulSoup(html, "html.parser")
    for tag in soup.select(JUNK):
        tag.decompose()
    return "\n".join(t.get_text(" ", strip=True)
                     for t in soup.select("p, h2, h3, li")
                     if t.get_text(strip=True))
\end{lstlisting}

Основной цикл (пачки случайных статей, параллельная загрузка):

\begin{lstlisting}[language=Python]
with ThreadPoolExecutor(workers) as pool:
    while count < limit:
        batch = [p for p in random_pages() if p[0] not in seen]
        seen.update(p[0] for p in batch)
        count += sum(pool.map(lambda p: download(p, folder), batch))
\end{lstlisting}

\pagebreak
\section{Выводы}
Выполнив первую лабораторную работу по курсу \enquote{Информационный поиск}, я
научился собирать корпус документов из открытого API. Основное, что я узнал
нового:

\begin{itemize}
    \item работа с MediaWiki API: выбор случайных статей, получение HTML,
    разбор ответов в формате JSON;
    \item выделение текста из HTML по CSS-селекторам и удаление служебных
    элементов (таблиц, сносок, навигации);
    \item обработка ограничений при загрузке большого объёма данных:
    повторные попытки, учёт кода 429 и \texttt{Retry-After}, возможность
    продолжить загрузку. На практике при слишком большом числе потоков (8) API
    начал отказывать, а при попытке обкачать новостной сайт соединение было
    заблокировано целиком, поэтому источник был заменён на Википедию;
    \item подсчёт статистики корпуса: «сырой» HTML в 9 раз больше полезного
    текста (69{,}1 КБ против 7{,}7 КБ на документ).
\end{itemize}

Корпус из 30027 статей будет использован в следующих лабораторных работах для
построения поискового робота, индекса и поиска. Встроенный поиск Википедии
показал недостатки --- неразрешённую многозначность, неочевидное ранжирование и
выдачу общих статей вместо точных; их предстоит учесть при построении
собственного поисковика.
